% TOP_PROBLEM: {"id": 2304001, "title": "Research Problems in Function Theory — Problem 4.1", "category_id": 8, "category": {"id": 8, "name": "analysis", "display_name": "Analysis", "description": "Limits, continuity, calculus, and function theory.", "slug": "analysis", "order_index": 8, "created_at": "2026-07-31T15:26:25.671Z"}, "status": "open"}
% FIRST_POSTED: null
% ENTRY_KIND: statement_only
% Imported from: batch13.tex; UnsolvedMath ID 2304001
\documentclass[11pt]{book}
\usepackage{fontspec}
\setmainfont{Latin Modern Roman}
\setsansfont{Latin Modern Sans}
\usepackage{amsmath,amssymb,mathtools}
\usepackage{unicode-math}
\setmathfont{Latin Modern Math}
\usepackage{xurl}
\usepackage[hidelinks]{hyperref}
\newcommand{\sref}[2]{\hyperlink{source:#1:#2}{[S#2]}}
\newcommand{\cataloguescope}[1]{\paragraph{Mathematical question.}#1}
\begin{document}
% UnsolvedMath ID 2304001; source catalogue code AMR-022-4001
\setcounter{section}{2304000}
\section{Research Problems in Function Theory — Problem 4.1}\label{2304001}
{\small\sffamily UnsolvedMath ID 2304001 · Analysis}
\subsection{Definitions and mathematical statement}

\paragraph{Shared notation.}$\mathbb N=\{1,2,\ldots\}$, and $\mathbb Z,\mathbb Q,\mathbb R,\mathbb C$ denote the integers, rationals, reals and complex numbers. Any other conventions are specified locally.


Write $\mathbb D=\{z\in\mathbb C:|z|<1\}$ and $\mathbb T=\partial\mathbb D$. Let $(z_j)_{j\ge1}$ be an infinite sequence of points on $\mathbb T$, with repetitions allowed. Define
\[A_n=\max_{|z|=1}\prod_{j=1}^n|z-z_j|,\qquad G_N=\left(\prod_{n=1}^NA_n\right)^{1/N}.
\]
The source's comparison sequence starts with $z_1=1$, $z_2=-1$ and, for each $k\ge1$, is extended by
\[z_{p+2^k}=z_p\exp\!\left(\frac{\pi i}{2^k}\right)\qquad(1\le p\le2^k).
\]
For this sequence $A_n\le n+1$, with equality exactly at $n=2^k-1$; values along $n=2^k$ remain bounded, so unboundedness is a statement about the upper limit rather than a pointwise limit.
\paragraph{Statement recorded in the supplied source.}
Must every infinite sequence satisfy $\limsup_{n\to\infty}A_n=+\infty$? Determine the sharp unavoidable growth of these upper-limit peaks and whether the displayed comparison sequence is extremal. The source's follow-up also asks for the possible and unavoidable growth of the geometric means $G_N$. \sref{2304001}{2}
\subsection{Short English statement}
Must the maximum distance products have arbitrarily large peaks, how small can their growth be, and is the stated dyadic construction optimal? Also study the geometric mean of all prefix maxima.
\subsection{Sources}
\begin{description}
\item[\hypertarget{source:2304001:1}{S1}] ULAM.AI and the UnsolvedMath Contributors, \emph{UnsolvedMath: A Curated Collection of Open Mathematics Problems}, record 2304001 (AMR-022-4001). CC BY 4.0. \url{https://huggingface.co/datasets/ulamai/UnsolvedMath}.
\item[\hypertarget{source:2304001:2}{S2}] Walter K. Hayman and Edward F. Lingham, \emph{Research Problems in Function Theory}, arXiv:1809.07200, 2018. Problem 4.1, its complete Update 4.1, and the relevant chapter notation. \url{https://arxiv.org/pdf/1809.07200}.
\end{description}
\label{end:2304001}
\end{document}
